% ============================================================================
%  VSEPR-SIM — Automatic Layering Report
%  Generated by the reporting pipeline with MiKTeX.
%
%  Compile: pdflatex layering_report.tex  (twice for TOC/refs)
%  Data:    layering_data.tex (auto-generated by generate_layering_report.py)
%
%  Reference: .github/copilot-instructions.md §2, §5, §9
% ============================================================================
\documentclass[11pt, letterpaper]{article}

% ── Geometry ────────────────────────────────────────────────────────────────
\usepackage[top=1.0in, bottom=1.0in, left=1.15in, right=1.15in,
            headheight=14pt]{geometry}

% ── Fonts / encoding ───────────────────────────────────────────────────────
\usepackage[T1]{fontenc}
\usepackage{lmodern}
\usepackage{microtype}

% ── Math ───────────────────────────────────────────────────────────────────
\usepackage{amsmath, amssymb}

% ── Tables ─────────────────────────────────────────────────────────────────
\usepackage{booktabs}
\usepackage{array}
\usepackage{longtable}

% ── Lists ──────────────────────────────────────────────────────────────────
\usepackage{enumitem}
\setlist{nosep, leftmargin=1.5em}

% ── Code listings ─────────────────────────────────────────────────────────
\usepackage{listings}
\usepackage{xcolor}

\definecolor{lstbg}{RGB}{248,248,248}
\definecolor{lstkw}{RGB}{0,0,180}
\definecolor{lstcm}{RGB}{80,120,80}

\lstdefinestyle{cppstyle}{
  language=C++, basicstyle=\ttfamily\footnotesize,
  backgroundcolor=\color{lstbg},
  keywordstyle=\color{lstkw}\bfseries,
  commentstyle=\color{lstcm}\itshape,
  frame=single, framesep=3pt, breaklines=true, keepspaces=true, tabsize=4,
  showstringspaces=false
}
\lstset{style=cppstyle}

% ── Colored boxes for layer summaries ─────────────────────────────────────
\usepackage{tcolorbox}
\tcbuselibrary{skins, breakable}

\definecolor{l1col}{RGB}{230,245,255}
\definecolor{l2col}{RGB}{255,245,220}
\definecolor{l3col}{RGB}{240,255,230}
\definecolor{l4col}{RGB}{255,230,240}
\definecolor{l5col}{RGB}{235,225,255}
\definecolor{crosscol}{RGB}{245,245,245}
\definecolor{hardcol}{RGB}{255,220,210}

% ── Headers / footers ─────────────────────────────────────────────────────
\usepackage{fancyhdr}
\pagestyle{fancy}
\fancyhf{}
\fancyhead[L]{\small\textit{VSEPR-SIM Research Platform}}
\fancyhead[R]{\small\textit{Automatic Layering Report}}
\fancyfoot[C]{\thepage}
\renewcommand{\headrulewidth}{0.4pt}

% ── Section formatting ─────────────────────────────────────────────────────
\usepackage{titlesec}
\titleformat{\section}{\Large\bfseries}{§\thesection}{1em}{}
\titleformat{\subsection}{\large\bfseries}{\thesubsection}{1em}{}
\titleformat{\subsubsection}{\normalsize\bfseries}{\thesubsubsection}{1em}{}

% ── Hyperlinks ─────────────────────────────────────────────────────────────
\usepackage{hyperref}
\hypersetup{colorlinks=true, linkcolor=blue!60!black, urlcolor=blue!70!black}

% ── Macros ─────────────────────────────────────────────────────────────────
\newcommand{\file}[1]{\texttt{#1}}
\newcommand{\type}[1]{\texttt{#1}}
\newcommand{\Ang}{\text{\AA}}
\newcommand{\kcalmol}{\text{kcal/mol}}
\newcommand{\lname}[1]{\textbf{L#1}}

% ── Import auto-generated data ─────────────────────────────────────────────
% Auto-generated by generate_layering_report.py
% Timestamp: 2026-06-28 22:55:39
% Source root: C:\R\VSPER-SIM

\newcommand{\ReportTimestamp}{2026-06-28 22:55:39}
\newcommand{\TotalFiles}{5398}
\newcommand{\TotalLines}{1,134,392}
\newcommand{\TotalHeaders}{2298}
\newcommand{\TotalSources}{3100}
\newcommand{\NumCMakeLibs}{0}
\newcommand{\NumBoundaryChecks}{5}
\newcommand{\NumGates}{2}
\newcommand{\NumReportModules}{9}

% --- Per-Layer Statistics ---
\newcommand{\LayerStatsTable}{
\begin{tabular}{lrrrrp{5.5cm}}
\toprule
\textbf{Layer} & \textbf{Files} & \textbf{Headers} & \textbf{Sources} & \textbf{Lines} & \textbf{Modules} \\
\midrule
L1 --- Paper Identity & 2 & 0 & 2 & 715 & \footnotesize Formation Priors \\
L2 --- Atomistic Kernel & 89 & 66 & 23 & 22,397 & \footnotesize Atomistic Core, Core, Integrators, Neighbor Lists, Potentials, Simulation \\
L3 --- Molecular 3D & 9 & 6 & 3 & 2,882 & \footnotesize Analysis, Atomistic Report, Validation \\
L4 --- Atomistic Bead (CG) & 47 & 47 & 0 & 11,615 & \footnotesize Atom-to-Bead Mapping, CG Core, CG Models, CG Reporting, Metals \\
L5 --- Macro / CAD / IO & 121 & 52 & 69 & 26,928 & \footnotesize API Facade, CLI, Demo, GUI, I/O, Rendering, Visualization \\
Cross-cutting & 652 & 329 & 323 & 166,907 & \footnotesize Applications, Examples, Include, Scripts, Tests, Tools \\
Unclassified & 4478 & 1798 & 2680 & 902,948 & \footnotesize Other \\
\midrule
\textbf{Total} & \textbf{5398} & \textbf{2298} & \textbf{3100} & \textbf{1,134,392} & \\
\bottomrule
\end{tabular}
}

% --- Layer Manifest ---
\newcommand{\LayerManifestTable}{
\begin{tabular}{llp{8cm}}
\toprule
\textbf{Layer} & \textbf{Status} & \textbf{Notes} \\
\midrule
Formation Priors & $\sim$ partial & \footnotesize Crystal presets complete. VSEPR seeds present. Learned topology prior absent. \\
Deterministic Physical Kernel & $\sim$ partial & \footnotesize LJ+PBC active. Bonded model implemented but not composed. Coulomb deferred. \\
Verification \& Benchmarking & $\checkmark$ complete & \footnotesize Phases 1-7 all pass. 194 checks. Session artifacts on disk. \\
Analysis \& Ranking & $\sim$ partial & \footnotesize Crystal metrics implemented. Structural ranking and learned prior absent. \\
Bridge \& Artifact Output & $\sim$ partial & \footnotesize EngineAdapter wires LoadXYZ/Save/Relax/MD/SinglePoint. EmitCrystal wired. \\
\bottomrule
\end{tabular}
}

% --- Boundary Checks ---
\newcommand{\BoundaryChecksTable}{
\begin{tabular}{lp{10cm}}
\toprule
\textbf{Check} & \textbf{Description} \\
\midrule
\texttt{B1} & \footnotesize Energy consistency check  Compare L2 pairwise energy (atomistic LJ+Coulomb) against L4  bead energy (role-weighted 3-channel). Record the mapping error. \\
\texttt{B2} & \footnotesize Charge conservation audit  Sum L2 charges for each bead group. Compare against L4 bead charge.  Flag non-conservation above threshold. \\
\texttt{B3} & \footnotesize Phase coherence check  Identify beads that straddle a phase boundary (atoms in the group  have inconsistent phase states). Flag and optionally split. \\
\texttt{B4} & \footnotesize Structural role re-evaluation  At the boundary, $\Sigma$\_i may need to be re-evaluated. An Fe atom  that was metallic in the ore body may become ionic after oxidation.  This unit re-classifies using the L2 charge and environment. \\
\texttt{B5} & \footnotesize Boundary mapping residual  Overall quality metric: how faithfully does L4 represent the L2  physics? High residual = the CG mapping is losing information here. \\
\bottomrule
\end{tabular}
}

% --- CMake Libraries ---
\newcommand{\CMakeLibsTable}{
\begin{tabular}{ll}
\toprule
\textbf{Target} & \textbf{Type} \\
\midrule
\bottomrule
\end{tabular}
}

% --- Report Modules ---
\newcommand{\ReportModulesTable}{
\begin{tabular}{lrl}
\toprule
\textbf{Module} & \textbf{Lines} & \textbf{File} \\
\midrule
Bead Fire Csv Export & 134 & \texttt{\footnotesize coarse\_grain/report/bead\_fire\_csv\_export.hpp} \\
Bead Fire Report & 283 & \texttt{\footnotesize coarse\_grain/report/bead\_fire\_report.hpp} \\
Excel Export & 222 & \texttt{\footnotesize coarse\_grain/report/excel\_export.hpp} \\
Formation Report & 549 & \texttt{\footnotesize coarse\_grain/report/formation\_report.hpp} \\
Mapping Report & 38 & \texttt{\footnotesize coarse\_grain/report/mapping\_report.hpp} \\
Metal Gas Study & 393 & \texttt{\footnotesize coarse\_grain/report/metal\_gas\_study.hpp} \\
Snapshot Graph & 374 & \texttt{\footnotesize coarse\_grain/report/snapshot\_graph.hpp} \\
Solidworks Export & 180 & \texttt{\footnotesize coarse\_grain/report/solidworks\_export.hpp} \\
Report Md & 7 & \texttt{\footnotesize atomistic/report/report\_md.hpp} \\
\bottomrule
\end{tabular}
}

% --- Gate Functions ---
\newcommand{\GateFunctions}{\texttt{l1\_gate}, \texttt{l2\_gate}}

% --- IKK Identity Metrics (D_rec / eta_ab per run, WO-75A) ---
\newcommand{\IKKSummaryTable}{
\textit{No \texttt{.identity.json} sidecar files found in \texttt{out/}.}
}



% ============================================================================
\begin{document}
% ============================================================================

% ── Title page ─────────────────────────────────────────────────────────────
\begin{titlepage}
  \centering
  \vspace*{2cm}
  {\Huge\bfseries VSEPR-SIM Research Platform\par}
  \vspace{0.8cm}
  {\Large\bfseries Automatic Layering Report\par}
  \vspace{0.4cm}
  \rule{0.7\textwidth}{0.6pt}\par
  \vspace{0.8cm}
  {\large\itshape
    Five-Layer Vertical Integration Stack\\[0.3em]
    Deterministic Atomistic Simulation Platform\\[0.3em]
    Auto-Generated from Source Tree Introspection\par}
  \vspace{1.2cm}
  \begin{tabular}{ll}
    \toprule
    \textbf{Metric} & \textbf{Value} \\
    \midrule
    Report generated & \ReportTimestamp \\
    Total source files & \TotalFiles \\
    Total non-empty lines & \TotalLines \\
    Header files & \TotalHeaders \\
    Source files & \TotalSources \\
    CMake library targets & \NumCMakeLibs \\
    Boundary checks & \NumBoundaryChecks \\
    Layer gates & \NumGates \\
    Report modules & \NumReportModules \\
    \bottomrule
  \end{tabular}
  \vfill
  {\small
    Implementation files:\quad
    \file{include/layer\_stack.hpp}\quad
    \file{atomistic/core/layer\_manifest.hpp}\quad
    \file{coarse\_grain/models/layer\_boundary.hpp}\par
    \vspace{0.3em}
    Compiled with MiKTeX \texttt{pdflatex}.
    Data generated by \file{reporting/generate\_layering\_report.py}.
  }
\end{titlepage}

\tableofcontents
\clearpage

% ============================================================================
\section{Architecture Overview}
\label{sec:overview}
% ============================================================================

The VSEPR-SIM platform is organized as a five-layer vertical integration
stack.  Each layer represents a different resolution of the same physical
object, from paper identity (formula and name) through atomistic physics,
molecular geometry, coarse-grained bead representation, to macro-scale
CAD geometry.

\subsection{The Five Layers}

\begin{tcolorbox}[colback=l1col, colframe=blue!40!black,
  title={\lname{1} --- Paper Identity}, fonttitle=\bfseries\small]
  Formula, class, name.  Static.  Never mutates.
  A candidate that fails L1 is never instantiated at L2 or above.
\end{tcolorbox}

\begin{tcolorbox}[colback=l2col, colframe=orange!60!black,
  title={\lname{2} --- Atomistic Kernel}, fonttitle=\bfseries\small]
  $Z$, $A$, $Q$, $\varepsilon$.  Physics and energy state.
  Environment-dependent.  The deterministic physical kernel lives here.
\end{tcolorbox}

\begin{tcolorbox}[colback=l3col, colframe=green!50!black,
  title={\lname{3} --- Molecular 3D}, fonttitle=\bfseries\small]
  Geometry, bonding, coordination, local environment.
  Built by FIRE minimization on the L2 state.
  Verification and validation benchmarks operate at this layer.
\end{tcolorbox}

\begin{tcolorbox}[colback=l4col, colframe=red!50!black,
  title={\lname{4} --- Atomistic Bead (CG)}, fonttitle=\bfseries\small]
  Coarse-grained beads, phase behaviour, bulk interactions.
  Anisotropic surface mapping and the three-channel interaction kernel
  live here.
\end{tcolorbox}

\begin{tcolorbox}[colback=l5col, colframe=violet!60!black,
  title={\lname{5} --- Macro / CAD / I/O}, fonttitle=\bfseries\small]
  Visualization, CLI, API facade, rendering, file I/O.
  The engineering and user-facing layer.
\end{tcolorbox}

\subsection{Vertical Integration Invariant}

For any particle $i$ and any layer $k \in \{2, 3, 4, 5\}$:
\begin{equation}
  L_k(i).Z = L_2(i).Z, \qquad L_k(i).A = L_2(i).A
  \label{eq:invariant}
\end{equation}

Nuclear identity is immutable across all layers under all non-nuclear
transformations.  Charge $Q_i$ and energy state $\varepsilon_i$ evolve
with chemical environment.

\subsection{Screening Funnel}

\begin{center}
\begin{tabular}{lrll}
\toprule
Layer & Typical survivors & Cost per candidate & Gate criterion \\
\midrule
L1 screen & $10^6$ & $\mu$s & Formula valid, class not excluded \\
L2 screen & $10^4$ & ms & $\varepsilon$ in range, $|Q| < Q_{\max}$ \\
L3 screen & $10^3$ & 0.1--1\,s & Geometry converged, $F_{\rm rms} < F_{\rm tol}$ \\
L4 screen & $10^2$ & 1--10\,s & Bead mapping residual $< \delta$, phase coherent \\
L5 report & $10^1$ & 10--100\,s & CAD geometry, integrity check \\
\bottomrule
\end{tabular}
\end{center}

% ============================================================================
\section{Layer Statistics}
\label{sec:stats}
% ============================================================================

The following table is auto-generated by walking the source tree and
classifying every C++ header and source file by its directory prefix.

\begin{table}[h!]
\centering
\caption{Per-layer file and line counts (auto-generated).}
\label{tab:layer-stats}
\small
\LayerStatsTable
\end{table}

% ============================================================================
\section{Layer Manifest}
\label{sec:manifest}
% ============================================================================

The authoritative layer status is declared in
\file{atomistic/core/layer\_manifest.hpp}.  This table is extracted
automatically from the \texttt{LAYER\_MANIFEST[]} array.

\begin{table}[h!]
\centering
\caption{Layer manifest status (from \texttt{layer\_manifest.hpp}).}
\label{tab:manifest}
\small
\LayerManifestTable
\end{table}

% ============================================================================
\section{L2 $\leftrightarrow$ L4 Boundary Checks}
\label{sec:boundary}
% ============================================================================

\begin{tcolorbox}[colback=hardcol, colframe=red!70!black,
  title={Hard Problem: The L2 $\leftrightarrow$ L4 Phase Transition Boundary},
  fonttitle=\bfseries\small, breakable]
This is where most simulations cheat or break.  Atomic behaviour
(discrete particle physics, quantum-derived $\varepsilon$, explicit
bonding) transitions into coarse-grained bulk behaviour (phase fields,
density order parameters, collective interactions).
\end{tcolorbox}

The boundary checks defined in \file{coarse\_grain/models/layer\_boundary.hpp}
are extracted below.

\begin{table}[h!]
\centering
\caption{L2 $\leftrightarrow$ L4 boundary checks (from \texttt{layer\_boundary.hpp}).}
\label{tab:boundary}
\small
\BoundaryChecksTable
\end{table}

\subsection{Boundary Mathematics}

Energy consistency (B1):
\begin{equation}
  \Delta E_{\text{boundary}} = |E_{\text{L4}}^{(\text{bead})} - E_{\text{L2}}^{(\text{atom})}|
  < \varepsilon_{\text{tol}}
\end{equation}

Charge conservation (B2):
\begin{equation}
  \Delta q = \left| q_{\text{L4}}^{(\text{bead})} - \sum_{j \in \text{group}} q_j^{(\text{L2})} \right|
  < q_{\text{tol}}
\end{equation}

Composite residual (B5):
\begin{equation}
  R = 0.5 \cdot E_{\text{term}} + 0.3 \cdot Q_{\text{term}} + 0.2 \cdot \phi_{\text{term}}
  < R_{\text{warn}}
\end{equation}

Structural role re-evaluation (B4):
\begin{equation}
  \Sigma_i^{(\text{boundary})} =
  \begin{cases}
    \text{IonicDominant}       & Z_i \in \text{transition metals} \;\wedge\; |Q_i| > 0.8 \\
    \text{Metallic}            & Z_i \in \text{transition metals} \;\wedge\; |Q_i| \leq 0.8 \\
    \text{IonicDominant}       & |Q_i| > 0.5 \\
    \text{DirectionalCovalent} & \text{otherwise}
  \end{cases}
\end{equation}

% ============================================================================
\section{Layer Gate Functions}
\label{sec:gates}
% ============================================================================

Gate functions are defined in \file{include/layer\_stack.hpp}.  Each gate
makes a binary accept/reject decision plus a score.  Accepted particles
propagate to the next layer; rejected particles are recorded with a reason.

\textbf{Detected gates:} \GateFunctions

The lazy evaluation strategy means that L3, L4, and L5 are represented as
\texttt{std::optional} in the \type{LayerParticle} struct.  No particle
is evaluated at layer $k$ until it has passed the gate at layer $k-1$:

\begin{lstlisting}
struct LayerParticle {
    L1_PaperIdentity                     l1;  // always populated
    L2_AtomisticState                    l2;  // populated after L1 gate
    std::optional<L3_MolecularGeometry>  l3;  // populated after L2 gate
    std::optional<L4_AtomisticBeadState> l4;  // populated after L3 gate
    std::optional<L5_MacroGeometry>      l5;  // populated after L4 gate
};
\end{lstlisting}

% ============================================================================
\section{CMake Build Targets}
\label{sec:cmake}
% ============================================================================

The following library targets are declared in the root
\file{CMakeLists.txt}.  Each target corresponds to a modular library
that can be linked independently.

\begin{table}[h!]
\centering
\caption{CMake library targets (auto-extracted).}
\label{tab:cmake}
\small
\CMakeLibsTable
\end{table}

% ============================================================================
\section{Report Modules}
\label{sec:report-modules}
% ============================================================================

The C++ report-generation modules produce structured outputs (Markdown,
CSV, XML, SolidWorks curves) that feed the downstream Python reporting
pipeline.

\begin{table}[h!]
\centering
\caption{C++ report modules (auto-discovered).}
\label{tab:report-modules}
\small
\ReportModulesTable
\end{table}

% ============================================================================
\section{Layer Dependency Graph}
\label{sec:deps}
% ============================================================================

The directed dependency graph of the five layers follows the vertical
integration invariant.  Higher layers depend on lower layers; identity
propagates upward.

\begin{center}
\begin{tabular}{lll}
\toprule
\textbf{From} & \textbf{To} & \textbf{Relationship} \\
\midrule
L1 & L2 & Formula $\to$ Atomistic state ($Z$, $A$, $\varepsilon$, $\sigma$) \\
L2 & L3 & Atomistic state $\to$ FIRE geometry optimization \\
L3 & L4 & Converged geometry $\to$ CG bead mapping \\
L2 $\leftrightarrow$ L4 & Boundary & Energy, charge, phase, role checks \\
L4 & L5 & Bead field $\to$ Macro body properties \\
L1--L5 & Report & All layers $\to$ Structured report output \\
\bottomrule
\end{tabular}
\end{center}

\subsection{FIRE Minimization (L2 $\to$ L3)}

\begin{equation}
  \mathbf{F}_i = -\nabla_{\mathbf{r}_i} U_{\text{total}}(\{\mathbf{r}\}, \{\mathbf{q}\})
\end{equation}

\begin{equation}
  \mathbf{r}_{i,t+1} = \mathbf{r}_{i,t} + \Delta t \, \mathbf{v}_{i,t}
\end{equation}

\begin{equation}
  \mathbf{v}_{i,t+1} = (1-\alpha)\mathbf{v}_{i,t}
    + \alpha\,\frac{\mathbf{F}_{i,t}}{\|\mathbf{F}_t\|}\,\|\mathbf{v}_{i,t}\|
\end{equation}

\begin{equation}
  \text{Converge if } \|\mathbf{F}\|_{\text{RMS}} < \varepsilon_F
  \;\lor\;
  \frac{|U_t - U_{t-1}|}{N} < \varepsilon_U
\end{equation}

\subsection{Anisotropic Surface Mapping (L4)}

The Fuzzy Ball Model at L4 carries a three-channel spherical harmonic
surface descriptor:
\begin{equation}
  S^{(k)}(\theta, \phi) = \sum_{\ell=0}^{\ell_{\max}}
                           \sum_{m=-\ell}^{+\ell}
                           c_{\ell m}^{(k)} Y_\ell^m(\theta, \phi),
  \quad k \in \{\text{steric}, \text{electrostatic}, \text{dispersion}\}
\end{equation}

% ============================================================================
\section{Output Properties}
\label{sec:output}
% ============================================================================

All report outputs satisfy the anti-black-box design philosophy:

\begin{itemize}
  \item \textbf{Explicit} --- every value maps to a specific computation.
  \item \textbf{Inspectable} --- Markdown, CSV, and LaTeX are human-readable.
  \item \textbf{Deterministic} --- same source tree yields identical report.
  \item \textbf{Modular} --- each report module operates independently.
  \item \textbf{Traceable} --- every metric links to a source file or struct.
\end{itemize}

\subsection{Regeneration}

This report can be regenerated at any time:

\begin{lstlisting}[language=bash]
# PowerShell (Windows / MiKTeX)
.\scripts\build_layering_report.ps1

# Manual steps
python reporting/generate_layering_report.py
cd reporting
pdflatex layering_report.tex
pdflatex layering_report.tex
\end{lstlisting}

% ============================================================================
\section*{Conclusion}
% ============================================================================

This layering report provides a machine-generated inventory of the
VSEPR-SIM five-layer vertical integration stack.  Every table, every
count, and every status field is derived from the current source tree
and C++ declarations --- not from hand-written prose.

The report is intended for:
\begin{itemize}
  \item Architectural auditing --- verify that layers are populated as expected
  \item Development tracking --- observe how file counts and line totals evolve
  \item Boundary monitoring --- ensure L2$\leftrightarrow$L4 checks remain comprehensive
  \item Research documentation --- embed in technical reports as an appendix
\end{itemize}

\vfill
\noindent\rule{\textwidth}{0.4pt}
\small\textit{Generated by VSEPR-SIM automated layering report pipeline.
Data: \file{reporting/generate\_layering\_report.py}.
Template: \file{reporting/layering\_report.tex}.
Compiled with MiKTeX \texttt{pdflatex}.}

\end{document}
